\documentclass[12pt]{article}
\usepackage[utf8]{inputenc}
\usepackage[T1]{fontenc}
\usepackage{hyperref}
\title{Ayi Language Resources for K-12 Curriculum Development}
\author{Compiled by Claude Sonnet 5.5 (claude-sonnet-5-5) for Andr\'as Kornai}
\date{}
\begin{document}
\maketitle

\section{Introduction}
Ayi (ISO 639-3 \texttt{ayq}) is a Sepik (Tama) language of Sandaun Province, Papua New Guinea. Sources give differing speaker counts (about 430 in the 2000 census; 900 in a Joshua Project figure based on Ethnologue). Very little material was found.

\section{Primary Educational Resources}
None found. No primers, textbooks or school programmes were found in the searches.

\section{Linguistic and Reference Materials}
\begin{itemize}
\item Ayi language, Wikipedia. \url{https://en.wikipedia.org/wiki/Ayi_language}
\item Glottolog entry for Ayi. \url{https://glottolog.org/resource/languoid/id/ayip1235}
\item Endangered Languages Project entry. \url{https://www.endangeredlanguages.com/lang/10698?hl=fr}
\end{itemize}

\section{Religious and Community Resources}
Joshua Project reports that the community has requested Bible translation since 1996 and lists the need as open, so no complete translation appears to exist. \url{https://joshuaproject.net/languages/ayq}

\section{Digital and Audio Resources}
None found.

\section{Curriculum Development Context}
The searches found no teaching or documentation materials in Ayi. A curriculum would start from new community-led documentation.

\end{document}